% Run the build script from any working directory: python3 thesis/current/build_thesis.py
% General review layout; replace with the institution's required template later.
\documentclass[UTF8,a4paper,12pt,fontset=fandol]{ctexrep}
\usepackage[margin=2.5cm]{geometry}
\usepackage{amsmath,amssymb,booktabs,graphicx,longtable,array}
\usepackage[hidelinks,unicode]{hyperref}
\hypersetup{pdftitle={基于执行反馈的小模型知识图谱问答方法研究及雷达领域应用},
  pdfsubject={电子信息专业硕士论文七章工作稿；AI参考与人工核验分别标注}}
\setlength{\emergencystretch}{2em}
\input{thesis/current/generated_tables.tex}
\input{thesis/current/data_tables.tex}
\input{thesis/current/domain_tables.tex}
\input{thesis/current/coverage_tables.tex}
\title{基于执行反馈的小模型知识图谱问答方法研究\\及雷达领域应用}
\author{电子信息专业硕士学位论文工作稿}
\date{2026年10月8日\quad 通用审阅版}
\begin{document}
\hypersetup{pageanchor=false}
\maketitle
\clearpage
\pagenumbering{Roman}
\hypersetup{pageanchor=true}
\input{thesis/current/00_abstract.tex}
\tableofcontents
\begingroup
\footnotesize
\listoftables
\endgroup
\clearpage
\pagenumbering{arabic}
\input{thesis/current/01_introduction.tex}
\input{thesis/current/02_background_related_work.tex}
\input{thesis/current/03_radar_data_construction.tex}
\input{thesis/current/04_execution_feedback_method.tex}
\input{thesis/current/05_public_experiments.tex}
\input{thesis/current/06_radar_application.tex}
\input{thesis/current/07_conclusion.tex}
\clearpage
\phantomsection
\addcontentsline{toc}{chapter}{参考文献}
\bibliographystyle{plain}
\bibliography{thesis/current/references}
\end{document}
